\chapter{Introduction}
\label{ch:introduction}

% The introduction follows the order in the writing rules (section 6):
%   1. general domain and method context;
%   2. current state of the field;
%   3. the research gap;
%   4. thesis statement and deliverables;
%   5. reader's guide.
% It motivates the problem in the domain, motivates the method challenge structurally,
% and explains why the thesis' framework matters. It is not a textbook review.

\section{Context}
\label{sec:ch1-context}

\todo{Domain and method context in two or three paragraphs.} The first paragraph
states the problem in the domain and why it matters. The second paragraph states the
method challenge and why existing approaches leave it open.

\section{State of the field}
\label{sec:ch1-field}

\todo{What is known, with citations.} Introduce each cited work on first mention and
reuse the key thereafter. Keep this section to the works the thesis builds on or
contrasts with.

\section{Research gap and thesis statement}
\label{sec:ch1-gap}

\todo{The gap in one paragraph; the thesis statement and the deliverables in another.}
State the research question so that each results chapter answers a part of it. Name the
deliverables in the same order as the chapters that produce them.

\section{Reader's guide}
\label{sec:ch1-guide}

\todo{One sentence per chapter.} Each sentence names the chapter and the question it
answers. A reader can then navigate to the part they need without reading the rest.
